\relax 
\GlossaryLetter{A}
\GlossaryEntry{绝对收敛}{Absolute Convergence}{无穷和中各项绝对值（复数项取模）的和有限；附录用它说明复指数展开可以按奇偶次幂分组。}
\GlossaryEntry{抽象向量空间}{Abstract Vector Space}{不限定元素具体形式，只要求加法与数乘满足向量空间公理的集合。}
\GlossaryEntry{激活函数}{Activation Function}{神经网络中对线性组合结果施加的函数，常逐分量作用并引入非线性；局部导数参与链式法则。}
\GlossaryEntry{主动旋转}{Active Rotation}{固定坐标基，对向量或物体施加旋转作用，用同一组基记录作用前后的坐标；转轴上的向量可以保持不变。}
\GlossaryEntry{加减消元法}{Addition Method of Elimination}{把一个方程的倍数加到另一个方程，以消去某个未知数的方法。}
\GlossaryEntry{映射的加法}{Addition of Maps}{对同定义域、同目标空间的映射逐点相加，$(S+T)(\boldsymbol x)=S(\boldsymbol x)+T(\boldsymbol x)$。}
\GlossaryEntry{相反向量}{Additive Inverse}{与原向量相加得到零向量的向量，记为 $-\boldsymbol v$。}
\GlossaryEntry{可加性}{Additivity}{映射保持向量加法的性质，即 $T(\boldsymbol x+\boldsymbol y)=T(\boldsymbol x)+T(\boldsymbol y)$。}
\GlossaryEntry{伴随矩阵}{Adjugate Matrix}{代数余子式数表的转置，记作 $\operatorname {adj}(A)$；其第 $i$ 列是第 $i$ 行对应的代数余子式向量，满足 $A\operatorname {adj}(A)=\det (A)E$。}
\GlossaryEntry{仿射映射}{Affine Map}{由线性作用与固定平移组成的映射 $\boldsymbol x\mapsto A\boldsymbol x+\boldsymbol a$；在原向量空间中，只有 $\boldsymbol a=\boldsymbol 0$ 时才是线性映射。}
\GlossaryEntry{仿射子空间}{Affine Subspace}{子空间的一次平移 $\boldsymbol p+W$；其维数为 $\dim W$，当 $\boldsymbol p\notin W$ 时不经过原点。}
\GlossaryEntry{交替最小二乘法}{Alternating Least Squares}{交替固定一组参数、对另一组参数求解最小二乘问题的方法；第3章以它说明推荐教学模型中的矩阵求解。}
\GlossaryEntry{交替多重线性型}{Alternating Multilinear Form}{交换任意两个输入位置使值变号的多重线性型；在实数范围，两个输入相同则值为零，归一化的交替$n$重线性型给出行列式。}
\GlossaryEntry{交替性}{Alternation}{本书实数情形中，交换两个输入使函数值变号的性质；与多重线性结合时，两输入相同则值为0。}
\GlossaryEntry{向量夹角}{Angle Between Vectors}{两个非零实向量的夹角，取值在 $[0,\pi ]$；余弦等于点积除以两向量长度之积。}
\GlossaryEntry{辐角}{Argument}{从正实轴转向非零复数所对应向量的有向角，相差 $2k\pi $ 的角表示同一方向。}
\GlossaryEntry{结合律}{Associative Law}{改变相邻运算的分组方式不改变结果；如 $(\boldsymbol u+\boldsymbol v)+\boldsymbol w=\boldsymbol u+(\boldsymbol v+\boldsymbol w)$。}
\GlossaryEntry{注意力机制}{Attention Mechanism}{根据查询与键计算权重，再对值向量加权求和的信息汇聚方法；第3章介绍性实例采用单头的简化图示。}
\GlossaryEntry{姿态矩阵}{Attitude Matrix}{常指机体系相对参考系的方向余弦矩阵；只改变参考基时按坐标转换链左乘。若记录同一个旋转作用并同时更换输入与输出的基，其表示矩阵才按相似关系改变。}
\GlossaryEntry{增广矩阵}{Augmented Matrix}{在系数矩阵右侧附上常数项列得到的矩阵，记为 $[A\mid \boldsymbol b]$。}
\GlossaryLetter{B}
\GlossaryEntry{回代}{Back Substitution}{从阶梯形的最后一个有效方程开始，逐行向上求出主元变量的方法。}
\GlossaryEntry{反向传播}{Backpropagation}{利用链式法则，从计算的输出端向前层依次传递梯度的方法；负责计算梯度，与使用梯度更新参数的优化步骤不同。}
\GlossaryEntry{主元变量}{Basic Variable}{系数部分的主元列所对应的未知数；有解时由自由变量确定。}
\GlossaryEntry{基}{Basis}{既线性无关又张成整个所讨论空间的向量组。}
\GlossaryEntry{基扩充定理}{Basis Extension Theorem}{有限维空间中的任意线性无关组，都能通过添加该空间中的向量扩充为一组基。}
\GlossaryEntry{双射}{Bijection}{同时为单射和满射的映射；每个目标元素恰好对应一个输入。}
\GlossaryEntry{分块对角矩阵}{Block Diagonal Matrix}{除主对角块之外各块均为零的矩阵；本书讨论的可逆性结论要求对角块为方阵。}
\GlossaryEntry{分块矩阵}{Block Matrix}{按若干固定行、列边界分成子矩阵的矩阵，用于记录分组输入与分组输出之间的作用。}
\GlossaryEntry{分块上三角矩阵}{Block Upper Triangular Matrix}{主对角块下方各块全为零的分块矩阵，可按组从后向前求解对应方程。}
\GlossaryLetter{C}
\GlossaryEntry{加法消去律}{Cancellation Law for Addition}{由 $\boldsymbol u+\boldsymbol w=\boldsymbol v+\boldsymbol w$ 可推出 $\boldsymbol u=\boldsymbol v$。}
\GlossaryEntry{笛卡尔积}{Cartesian Product}{集合按顺序组合形成有序组的集合；$V^m$的每个元素是来自$V$的$m$个向量组成的有序组。}
\GlossaryEntry{柯西--施瓦茨不等式}{Cauchy-Schwarz Inequality}{实向量点积的绝对值不超过两向量长度之积；保证非零向量夹角余弦公式的取值在 $[-1,1]$ 内，第7章给出证明。}
\GlossaryEntry{中心}{Center}{附近非平衡状态沿闭合轨线绕行的典型矢量场形态；本书用二维圆周运动模型建立直观。}
\GlossaryEntry{链式法则}{Chain Rule}{复合函数的局部变化率由各层局部映射复合得到；多元情形按正确取值位置将雅可比矩阵依次相乘。}
\GlossaryEntry{基变换}{Change of Basis}{改变描述向量所用的有序基，从而改变坐标表示；向量本身不变。第4章选读讨论坐标转换。}
\GlossaryEntry{过渡矩阵}{Change-of-Coordinates Matrix}{把同一向量在一组有序基下的坐标转换到另一组基下的坐标的可逆矩阵；必须同时指明输入基与输出基。}
\GlossaryEntry{对运算封闭}{Closure under an Operation}{集合中允许的元素按指定运算组合后，结果仍属于该集合。}
\GlossaryEntry{陪域}{Codomain}{定义映射时预先指定的目标集合；对 $T:\mathbb R^n\to \mathbb R^m$，陪域为 $\mathbb R^m$，其中的元素不一定都被取到。}
\GlossaryEntry{系数}{Coefficient}{在给定表达式中乘在未知数或向量前的数。}
\GlossaryEntry{系数矩阵}{Coefficient Matrix}{按固定未知数次序逐行收集线性方程组系数形成的数表。}
\GlossaryEntry{代数余子式}{Cofactor}{余子式乘位置符号所得的数，$A_{ij}=(-1)^{i+j}M_{ij}$；本章矩阵元素另记作小写 $a_{ij}$。}
\GlossaryEntry{代数余子式向量}{Cofactor Vector}{由某一行的全部代数余子式组成的向量，记作 $\boldsymbol c_i$；垂直于其余各行，与该行的点积为行列式。}
\GlossaryEntry{冷启动}{Cold Start}{新用户或新内容缺少交互记录、相应参数尚未充分训练的推荐问题；不能简单等同于向量不在某个子空间中。}
\GlossaryEntry{共线}{Collinear}{向量起点统一后位于同一直线上；两个非零向量共线当且仅当一个是另一个的数倍。}
\GlossaryEntry{列秩}{Column Rank}{矩阵列向量张成空间的维数；等于矩阵的秩，也等于行秩。}
\GlossaryEntry{列空间}{Column Space}{矩阵各列向量张成的空间，即全部可能输出组成的集合；$m\times n$ 矩阵的列空间是 $\mathbb R^m$ 的子空间。第3.5节从转置引入。}
\GlossaryEntry{列向量}{Column Vector}{将分量按固定次序竖排得到的向量形式。}
\GlossaryEntry{交换图}{Commutative Diagram}{用箭头表示映射，且任意两条起终点相同的相配路径具有相同的合成作用；本书用它说明换基与矩阵表示的关系。}
\GlossaryEntry{交换律}{Commutative Law}{交换运算对象的次序不改变结果；如 $\boldsymbol u+\boldsymbol v=\boldsymbol v+\boldsymbol u$。}
\GlossaryEntry{共轭复数}{Complex Conjugate}{与 $z=a+bi$ 对应的 $\overline z=a-bi$，在复平面上关于实轴镜像对称。}
\GlossaryEntry{复指数}{Complex Exponential}{对实数 $a,\theta $，$e^{a+i\theta }=e^a(\cos \theta +i\sin \theta )$；模为 $e^a$，辐角可取 $\theta $。}
\GlossaryEntry{复数}{Complex Number}{形如 $a+bi$ 的数，其中 $a,b\in \mathbb R$，$i^2=-1$。}
\GlossaryEntry{复平面}{Complex Plane}{将复数 $a+bi$ 对应到点 $(a,b)$ 的坐标平面，横轴为实轴，纵轴为虚轴。}
\GlossaryEntry{复向量空间}{Complex Vector Space}{以复数作为数乘系数的向量空间；讨论线性关系时必须允许复系数。}
\GlossaryEntry{分量}{Component}{有序数组中某个固定位置的数，位置次序是向量表示的一部分。}
\GlossaryEntry{复合}{Composition}{把一个映射的输出作为另一个映射的输入；$(S\circ T)(\boldsymbol x)=S(T(\boldsymbol x))$，先作右侧映射。}
\GlossaryEntry{保守力}{Conservative Force}{可由单值势能函数的负梯度表示的力；力的方向指向势能下降的方向。}
\GlossaryEntry{有解}{Consistent}{线性方程组至少有一个公共解，也称相容。}
\GlossaryEntry{常数项}{Constant Term}{线性方程中不含未知数的已知量；本书通常将其放在等号右端。}
\GlossaryEntry{内容空间}{Content Space}{教学推荐案例中由给定视频特征向量张成的子空间；不是有限视频集合本身。}
\GlossaryEntry{坐标向量}{Coordinate Vector}{将向量关于有序基的表示系数按次序排成的列向量，记为 $[\boldsymbol v]_{\mathcal B}$。}
\GlossaryEntry{坐标}{Coordinates}{向量关于一组有序基的唯一线性表示系数，随所选基及其顺序而定。}
\GlossaryEntry{共面}{Coplanar}{向量起点统一后位于同一个平面内；三个空间向量共面时线性相关。}
\GlossaryEntry{余弦相似度}{Cosine Similarity}{两个非零实向量的标准内积除以各自长度的乘积，用于比较方向的接近程度。}
\GlossaryEntry{叉乘}{Cross Product}{两支三维向量所确定的法向量，长度为平行四边形面积，方向按右手法则；共线时为零向量。}
\GlossaryEntry{旋度}{Curl}{三维向量场横向偏导数按叉乘结构成对作差得到的向量，描述局部旋转倾向；刚体旋转场的旋度为角速度向量的两倍。}
\GlossaryLetter{D}
\GlossaryEntry{棣莫弗定理}{De Moivre's Theorem}{对整数 $n$，$(\cos \theta +i\sin \theta )^n=\cos (n\theta )+i\sin (n\theta )$。}
\GlossaryEntry{自由度}{Degrees of Freedom}{表示全部解所需的独立自由参数个数；线性方程组有解时等于相应零空间的维数。}
\GlossaryEntry{导数}{Derivative}{函数相对于自变量的瞬时变化率，记作 $f'$；求导保持函数的加法与实数数乘。}
\GlossaryEntry{互相关矩阵}{Detector Correlation Matrix}{本书在特征探测器语境中指矩阵 $A^{\mathrm T}A$；其第 $(i,j)$ 个元素为列向量的点积 $\boldsymbol a_i\cdot \boldsymbol a_j$，对角线记录自身增益，非对角线记录重叠响应；这里未作归一化，不是统计学的相关系数矩阵。}
\GlossaryEntry{行列式}{Determinant}{与方阵对应的标量，绝对值是线性变换的体积缩放倍数，符号记录方向是否翻转；为零时方阵不可逆。第6章展开。}
\GlossaryEntry{对角矩阵}{Diagonal Matrix}{主对角线以外的元素全为零的方阵，对应各坐标分量分别缩放。}
\GlossaryEntry{矩阵对角化}{Diagonalization}{选择适当的基，使某个线性变换的矩阵表示成为对角矩阵；并非所有方阵都能对角化。第9章展开。}
\GlossaryEntry{维数}{Dimension}{有限维空间中任意一组基所含的向量个数，记为 $\dim V$；零子空间的维数为0。}
\GlossaryEntry{维数公式}{Dimension Formula}{有限维子空间满足 $\dim (U+V)=\dim U+\dim V-\dim (U\cap V)$。}
\GlossaryEntry{直和}{Direct Sum}{当 $U\cap V=\{\boldsymbol 0\}$ 时的子空间和，记作 $U\oplus V$；等价地，每个和向量的分解唯一。}
\GlossaryEntry{有向角}{Directed Angle}{同时记录转动大小与方向的角；本附录约定逆时针为正、顺时针为负。}
\GlossaryEntry{方向余弦矩阵}{Direction Cosine Matrix}{两组单位且两两垂直的基之间的坐标过渡矩阵；第 $(i,j)$ 个元素为旧基第 $i$ 个向量与新基第 $j$ 个向量的夹角余弦，须注明换算方向。}
\GlossaryEntry{方向导数}{Directional Derivative}{沿指定单位方向的瞬时变化率；对本附录的平滑函数，等于梯度与该单位方向的点积，是带符号的投影长度。}
\GlossaryEntry{离散余弦变换}{Discrete Cosine Transform}{把有限个采样值改写为不同余弦模式的系数；本书使用规范化的二型DCT，完整保留系数时可以逆变换恢复输入。}
\GlossaryEntry{离散线性动力系统}{Discrete-Time Linear Dynamical System}{按离散时刻用线性规则更新状态的系统，如 $\boldsymbol x_{k+1}=F\boldsymbol x_k+\boldsymbol u_k$；解是一串状态记录。}
\GlossaryEntry{分配律}{Distributive Law}{数乘对向量加法与标量加法分别满足分配规则。}
\GlossaryEntry{散度}{Divergence}{向量场各分量对相同坐标方向的偏导数之和，即场的雅可比矩阵的迹；表示局部单位体积的净流出率。}
\GlossaryEntry{定义域}{Domain}{映射允许的全部输入组成的集合；对 $T:\mathbb R^n\to \mathbb R^m$，定义域为 $\mathbb R^n$。}
\GlossaryEntry{数量积（点积）}{Dot Product}{实向量对应分量乘积之和；在二维、三维中等于两向量长度之积乘夹角余弦，零向量与任意向量的点积为0。}
\GlossaryEntry{二重积分}{Double Integral}{对平面区域上小面积的函数贡献进行累积的积分；可表示曲面下的体积或由面密度决定的总质量。}
\GlossaryEntry{对偶基}{Dual Basis}{与一组基对应的坐标提取泛函所组成的基；标准基的对偶基由 $\alpha _j(\boldsymbol x)=x_j$ 这些坐标尺子组成，每个线性评分器由它们唯一线性组合得到。}
\GlossaryEntry{对偶映射}{Dual Map}{线性映射 $T:V\to W$ 所诱导的反向评价传递 $T':W'\to V'$，定义为 $T'(\varphi )=\varphi \circ T$；在标准对偶基下由转置矩阵表示，不是逆映射。}
\GlossaryEntry{对偶空间}{Dual Space}{实向量空间 $V$ 上全体线性泛函按逐点加法与数乘组成的向量空间，记为 $V'$；有限维时与 $V$ 维数相等，可理解为所有线性测量尺子的集合。}
\GlossaryLetter{E}
\GlossaryEntry{特征空间}{Eigenspace}{给定特征值对应的全部特征向量连同零向量构成的子空间；第8章展开。}
\GlossaryEntry{特征值}{Eigenvalue}{线性变换在某个非零向量方向上表现为数乘时的缩放系数；第8章展开。}
\GlossaryEntry{初等列变换}{Elementary Column Operations}{交换两列、非零倍乘一列、把一列的倍数加到另一列，分别对应右乘合适的初等矩阵。}
\GlossaryEntry{同解初等变换}{Elementary Equation Operations}{交换方程、非零倍乘一个方程、把一方程的倍数加到另一方程的操作，均保持解集不变。}
\GlossaryEntry{初等矩阵}{Elementary Matrix}{由单位矩阵作一次初等行变换得到的矩阵；左乘它等价于对原矩阵作同一次行变换。}
\GlossaryEntry{初等行变换}{Elementary Row Operations}{交换两行、非零倍乘一行、把一行的倍数加到另一行的三种矩阵操作。}
\GlossaryEntry{消元乘数}{Elimination Multiplier}{也称消元倍数；行消元 $r_i\leftarrow r_i-\ell _{i,k}r_k$ 中的系数 $\ell _{i,k}$；无换行的标准LU消元将它记录在 $L$ 的第 $(i,k)$ 位置。}
\GlossaryEntry{向量嵌入}{Embedding}{把对象或特征表示为供模型计算的向量，其分量不必对应人为命名的现实属性。}
\GlossaryEntry{嵌入表}{Embedding Table}{按对象或特征ID索引向量参数的数据组织方式；数学上可将有限表理解为每个ID对应一行的矩阵。}
\GlossaryEntry{平衡状态}{Equilibrium State}{更新后或随时间不改变的状态；离散系统满足 $(E-F)\boldsymbol x_*=\boldsymbol u$，齐次连续系统满足 $A\boldsymbol y_*=\boldsymbol 0$。}
\GlossaryEntry{同解方程组}{Equivalent Linear Systems}{解集相同的两个线性方程组。}
\GlossaryEntry{等价向量组}{Equivalent Sets of Vectors}{能够相互线性表示的向量组；其张成空间相同，但向量个数可以不同。}
\GlossaryEntry{模长}{Euclidean Length}{实向量各分量平方和的平方根，等于向量与自身点积的平方根；也称向量的长度。}
\GlossaryEntry{欧拉角}{Euler Angles}{以约定轴序的三次旋转描述姿态的角度参数；须说明转轴固定还是随动。某些姿态可由不同角度组表示，本书也以此广义称呼三个不同轴的角度参数。}
\GlossaryEntry{欧拉公式}{Euler's Formula}{对实数 $\theta $，$e^{i\theta }=\cos \theta +i\sin \theta $，将单位圆上的旋转与复指数联系起来。}
\GlossaryEntry{指数形式}{Exponential Form}{用 $z=re^{i\theta }$ 表示非零复数的形式，其中 $r=|z|$，$\theta $ 是其一个辐角。}
\GlossaryLetter{F}
\GlossaryEntry{特征向量（数据表示）}{Feature Vector}{按约定顺序记录对象各项特征的向量；此处不是后续特征值理论中的特征向量。}
\GlossaryEntry{域}{Field}{含不同的0与1，支持加减乘除（除数非零）并满足相应基本运算律的数系。}
\GlossaryEntry{复数域}{Field of Complex Numbers}{复数全体 $\mathbb C$ 连同通常的加法和乘法，构成一个域。}
\GlossaryEntry{有理数域}{Field of Rational Numbers}{有理数全体 $\mathbb Q$ 连同通常的加法和乘法，构成一个域。}
\GlossaryEntry{实数域}{Field of Real Numbers}{实数全体 $\mathbb R$ 连同通常的加法和乘法；本书主线所用的标量来源。}
\GlossaryEntry{有限域}{Finite Field}{只含有限个元素的域；本书仅作拓展提及。}
\GlossaryEntry{有限维空间}{Finite-Dimensional Space}{具有有限基的向量空间，零空间也属于有限维空间。}
\GlossaryEntry{形式行列式}{Formal Determinant}{含标准基向量的一行仅作为按行展开的记号，用于把各个标量余子式组成一个向量；不是普通的标量行列式，其定向符号取决于标准基行所在的位置。}
\GlossaryEntry{向前消元}{Forward Elimination}{逐列选择主元并消去其下方元素，将方程组或增广矩阵化成阶梯形的过程。}
\GlossaryEntry{前代}{Forward Substitution}{从下三角方程组的第1行开始，自上而下依次求出未知量的方法。}
\GlossaryEntry{四个基本子空间}{Four Fundamental Subspaces}{矩阵的列空间、行空间、零空间与左零空间；行空间和零空间位于输入空间，列空间和左零空间位于输出空间。第3.5节铺垫，后续研究维数与正交关系。}
\GlossaryEntry{自由变量}{Free Variable}{系数矩阵非主元列对应的未知数；有解时可独立指定其值，再求主元变量。}
\GlossaryEntry{自由向量}{Free Vector}{只由大小和方向确定、允许平移而不改变自身的几何向量。}
\GlossaryEntry{列满秩}{Full Column Rank}{矩阵的列向量全部线性无关，即其列向量组的秩等于列数。}
\GlossaryEntry{行满秩}{Full Row Rank}{矩阵的行向量全部线性无关，即其行向量组的秩等于行数。}
\GlossaryEntry{函数}{Function}{为定义域中的每个元素指定唯一输出的对应规则；函数与映射是同一种概念，输入和输出可以是数，也可以是向量等对象。}
\GlossaryEntry{基础解系}{Fundamental Set of Solutions}{齐次线性方程组零空间的一组基；每个齐次解均可唯一表示为这些基向量的线性组合，向量个数为未知数个数减系数矩阵的秩。}
\GlossaryEntry{平面向量基本定理}{Fundamental Theorem of Plane Vectors}{同一平面内任意向量可唯一表示为两个不共线向量的线性组合。}
\GlossaryEntry{空间向量基本定理}{Fundamental Theorem of Space Vectors}{空间中任意向量可唯一表示为三个不共面向量的线性组合。}
\GlossaryLetter{G}
\GlossaryEntry{高斯--若尔当消元法}{Gauss-Jordan Elimination}{继续把主元化为1并消去主元列其余元素，得到行最简形的方法。}
\GlossaryEntry{高斯消元法}{Gaussian Elimination}{通过同解初等变换化为行阶梯形，再指定自由变量并回代求解的方法。}
\GlossaryEntry{通解}{General Solution}{以任意参数或自由的齐次解完整表示全部解的表达式；不是一个特解或有限几个解。}
\GlossaryEntry{广义叉乘}{Generalized Cross Product}{定向的欧氏$n$维空间中，由$n-1$条有序边确定的法向量；长度等于其底面体积，最后接上该法向量构成正定向，边线性相关时为零。}
\GlossaryEntry{梯度}{Gradient}{把标量函数对各输入坐标的偏导数按顺序排成的列向量；对平滑函数，非零梯度指向每单位距离上升最快的方向。}
\GlossaryEntry{梯度下降}{Gradient Descent}{按负梯度方向小步调整参数的方法；梯度非零且步长足够小时局部下降，不自动保证达到全局最低点。}
\GlossaryLetter{H}
\GlossaryEntry{齐次性}{Homogeneity}{映射保持数乘的性质，即对所有实数 $c$ 和允许的向量 $\boldsymbol x$ 有 $T(c\boldsymbol x)=cT(\boldsymbol x)$。}
\GlossaryEntry{齐次坐标}{Homogeneous Coordinates}{用增加一个分量的方式表示几何点；本章把二维点记为 $(x,y,1)^{\mathrm T}$，以三阶矩阵统一表示平移与线性作用。}
\GlossaryEntry{齐次线性方程组}{Homogeneous Linear System}{常数项全为0的线性方程组 $A\boldsymbol x=\boldsymbol 0$，必有零解。}
\GlossaryEntry{齐次递推}{Homogeneous Recurrence}{没有外部输入的线性状态更新，如 $\boldsymbol x_{k+1}=F\boldsymbol x_k$；全部解序列对逐项加法和数乘封闭。}
\GlossaryEntry{齐次解空间}{Homogeneous Solution Space}{齐次线性方程全部解在相应加法与数乘下构成的向量空间；动态系统中，其元素是解序列或解函数，不能与平衡状态的零空间混同。}
\GlossaryEntry{超平面}{Hyperplane}{$\mathbb R^n$ 中维数为 $n-1$ 的仿射子空间；可由一条系数不全为零的线性方程描述。}
\GlossaryLetter{I}
\GlossaryEntry{恒等映射}{Identity Map}{把每个输入原样输出的映射，$I(\boldsymbol x)=\boldsymbol x$，其标准基矩阵为单位矩阵。}
\GlossaryEntry{单位矩阵}{Identity Matrix}{主对角线全为1、其余元素全为0的方阵，本书统一记为 $E_n$ 或 $E$。}
\GlossaryEntry{像}{Image}{给定输入在映射下对应的输出；$\boldsymbol x$ 在 $T$ 下的像记为 $T(\boldsymbol x)$。}
\GlossaryEntry{虚轴}{Imaginary Axis}{复平面中实部为0的点组成的坐标轴。}
\GlossaryEntry{虚部}{Imaginary Part}{复数 $a+bi$ 中的实数 $b$，不是 $bi$。}
\GlossaryEntry{虚数单位}{Imaginary Unit}{满足 $i^2=-1$ 的数 $i$；它的连续整数次幂以4为周期重复。}
\GlossaryEntry{关联矩阵}{Incidence Matrix}{记录网络节点与线路连接的矩阵；本书取节点为行、定向线路为列，每列在流出节点记$+1$、流入节点记$-1$，其余为$0$。}
\GlossaryEntry{无解}{Inconsistent}{线性方程组没有公共解，也称不相容；消元后可见 $0=c$、$c\ne 0$ 的矛盾。}
\GlossaryEntry{无限维空间}{Infinite-Dimensional Space}{不存在有限基的向量空间；例如实连续函数空间 $C[a,b]$，其中 $a<b$。}
\GlossaryEntry{无穷多解}{Infinitely Many Solutions}{实线性方程组有解且存在自由变量时的情形，对应无穷多个不同的解向量。}
\GlossaryEntry{单射}{Injection}{不同输入具有不同输出的映射；等价地，每个输出至多对应一个输入。}
\GlossaryEntry{内积}{Inner Product}{用于定义长度、夹角及正交性的一种向量配对运算；第7章正式定义。}
\GlossaryEntry{内积空间}{Inner Product Space}{配有指定内积的向量空间，使长度与正交性有准确含义；第7章展开。}
\GlossaryEntry{子空间的交}{Intersection of Subspaces}{同时属于两个子空间的向量全体，仍是子空间，记为 $U\cap V$。}
\GlossaryEntry{反函数}{Inverse Function}{在指定定义域与目标集合之间，双射函数把输出唯一还原为输入的函数；它与原函数两个方向的复合均为相应集合的恒等映射。}
\GlossaryEntry{逆映射}{Inverse Map}{双射把每个输出反查为唯一输入的映射；双射线性映射的逆仍是线性映射。}
\GlossaryEntry{逆矩阵}{Inverse Matrix}{满足 $AB=BA=E$ 的方阵 $B$，记为 $A^{-1}$，表示撤销 $A$ 对应线性作用的映射。}
\GlossaryEntry{逆序数}{Inversion Number}{排列中前项大于后项的数对个数；行列式项的符号为 $(-1)^{N(\sigma )}$。}
\GlossaryEntry{可逆矩阵}{Invertible Matrix}{存在双侧逆的方阵；等价于对应线性映射为双射。}
\GlossaryLetter{J}
\GlossaryEntry{雅可比行列式}{Jacobian Determinant}{输入与输出维数相同时，雅可比矩阵的行列式；绝对值表示局部的面积或体积缩放倍率，正负记录方向。}
\GlossaryEntry{雅可比矩阵}{Jacobian Matrix}{按输出分量为行、输入坐标为列排列偏导数的矩阵；在固定一点，把输入小变化映到输出的一阶变化。}
\GlossaryLetter{K}
\GlossaryEntry{键}{Key}{注意力计算中与查询配对、用来计算关联分数的向量，常按行组成矩阵 $K$。}
\GlossaryLetter{L}
\GlossaryEntry{拉普拉斯展开}{Laplace Expansion}{沿固定的一行或一列，将元素与对应代数余子式的乘积求和，得到整个行列式的计算公式。}
\GlossaryEntry{隐因子}{Latent Factor}{模型用于表示潜在联系的向量分量；本章以少量隐因子展示用户与内容的匹配。}
\GlossaryEntry{余弦定理}{Law of Cosines}{三角形一边的平方等于另两边的平方和减去这两边与其夹角余弦乘积的两倍；将勾股定理推广到任意三角形。}
\GlossaryEntry{学习率}{Learning Rate}{梯度下降中控制每次参数更新步长的正数；太大的步长可能使损失增加。}
\GlossaryEntry{最小二乘法}{Least Squares}{通过使误差平方和尽可能小来确定参数的方法；第3章先应用其线性方程形式，第5.7节用正交直和说明最优近似与解的结构。}
\GlossaryEntry{最小二乘解}{Least-Squares Solution}{使 $\|\boldsymbol b-A\boldsymbol x\|^2$ 最小的输入；等价于正规方程的解，最优输出唯一，输入是否唯一取决于 $A$ 的零空间。}
\GlossaryEntry{左零空间}{Left Null Space}{转置矩阵的零空间，即满足 $A^{\mathrm T}\boldsymbol y=\boldsymbol 0$ 的全部向量；等价地有 $\boldsymbol y^{\mathrm T}A=\boldsymbol 0^{\mathrm T}$。若 $A$ 为 $m\times n$，它是 $\mathbb R^m$ 的子空间。}
\GlossaryEntry{等高线}{Level Curve}{二元标量函数中具有相同函数值的输入点组成的曲线；非零梯度垂直于经过该点的等高线切向。}
\GlossaryEntry{线性代数}{Linear Algebra}{研究向量的线性关系、向量空间与线性映射及其代数表示的数学分支。}
\GlossaryEntry{线性组合}{Linear Combination}{对给定向量分别作数乘后相加得到的向量。}
\GlossaryEntry{线性方程}{Linear Equation}{形如 $a_1x_1+\cdots +a_nx_n=b$ 的方程，未知数只以一次项出现。}
\GlossaryEntry{线性泛函}{Linear Functional}{从实向量空间到 $\mathbb R$ 的线性评分规则，如坐标提取、固定权重点积和多项式在固定点的取值；在 $\mathbb R^m$ 中可唯一写成 $\varphi _{\boldsymbol y}(\boldsymbol b)=\boldsymbol y\cdot \boldsymbol b$。}
\GlossaryEntry{线性映射}{Linear Map}{保持向量加法和数乘的映射。在线性代数中也称线性变换；固定定义域与目标空间的标准基后，由各个基向量的像唯一确定其矩阵。}
\GlossaryEntry{线性运算}{Linear Operations}{向量加法与数乘的统称。}
\GlossaryEntry{线性算子}{Linear Operator}{从一个向量空间到它自身的线性映射；例如，次数不超过 $n$ 的多项式空间上的求导映射 $D(f)=f'$。}
\GlossaryEntry{线性表示}{Linear Representation}{给定向量可写成某向量组的线性组合这一关系。}
\GlossaryEntry{线性结构}{Linear Structure}{由向量加法、数乘及其线性表示关系构成的结构；其研究不要求预先给定长度与夹角。}
\GlossaryEntry{线性方程组}{Linear System}{要求同一组未知数同时满足的若干线性方程。}
\GlossaryEntry{线性微分方程组}{Linear System of Differential Equations}{以状态分量的导数为未知变化率的线性关系，如 $\dot {\boldsymbol y}=A\boldsymbol y+\boldsymbol g(t)$；其解是随时间变化的函数。}
\GlossaryEntry{线性变换}{Linear Transformation}{保持加法和数乘的向量映射，即 $T(\boldsymbol u+\boldsymbol v)=T\boldsymbol u+T\boldsymbol v$、$T(c\boldsymbol u)=cT\boldsymbol u$。}
\GlossaryEntry{线性相关}{Linearly Dependent}{存在不全为零的系数，使给定向量组的线性组合等于零向量。}
\GlossaryEntry{线性无关}{Linearly Independent}{只有全部系数为零，给定向量组的线性组合才等于零向量。}
\GlossaryEntry{荷载工况}{Load Case}{结构分析中一组指定的外部荷载情形；结构参数固定时，不同工况通常对应不同右端项。}
\GlossaryEntry{局部线性化}{Local Linearization}{在固定位置，用一个线性映射近似输入小变化引起的输出变化；加上起点输出后的近似通常是仿射映射，不是把原函数变成全局线性函数。}
\GlossaryEntry{损失函数}{Loss Function}{用一个标量衡量模型预测与目标的偏差或其他训练目标的函数；把模型参数视为输入时，是一个多元标量函数。}
\GlossaryEntry{下三角矩阵}{Lower Triangular Matrix}{主对角线上方全为零的方阵。}
\GlossaryEntry{LU分解}{LU Factorization}{将矩阵写成 $A=LU$，其中 $L$ 为单位下三角矩阵，$U$ 为上三角矩阵，用来保存并复用无换行消元。}
\GlossaryLetter{M}
\GlossaryEntry{主对角线}{Main Diagonal}{矩阵中行下标与列下标相等的位置，即 $A_{1,1},A_{2,2},\ldots $ 所在的对角线。}
\GlossaryEntry{映射}{Mapping}{给每个输入指定唯一输出的对应规则。}
\GlossaryEntry{矩阵}{Matrix}{按固定行列次序排列的矩形数表，可组织方程系数或记录线性映射的信息。}
\GlossaryEntry{矩阵加法}{Matrix Addition}{同型矩阵按对应位置元素相加的运算，表示相应线性映射输出的逐点叠加。}
\GlossaryEntry{矩阵乘法}{Matrix Multiplication}{表示线性映射复合的运算；相邻尺寸匹配时，乘积元素由左矩阵的一行与右矩阵的一列对应相乘求和。}
\GlossaryEntry{矩阵的幂}{Matrix Power}{方阵与自身反复相乘，表示同一映射反复复合；约定 $A^0=E$。}
\GlossaryEntry{矩阵表示}{Matrix Representation}{选定两端的有序基后，按列记录定义域各基向量的像在目标空间基下的坐标所得的矩阵；本书第3章使用标准基，记为 $M(T)$。}
\GlossaryEntry{矩阵数乘}{Matrix Scalar Multiplication}{将矩阵每个元素同乘一个标量，表示把线性映射的所有输出作相同倍数的缩放。}
\GlossaryEntry{极大线性无关组}{Maximal Linearly Independent Subset}{从原向量组中选出的无关部分组，并能线性表示原组的全部向量。}
\GlossaryEntry{度量结构}{Metric Structure}{用于刻画长度、夹角等几何量的附加结构；本书第0章从实向量的标准点积出发建立直观。}
\GlossaryEntry{余子式}{Minor}{删去指定元素所在的一行和一列，保留其余行列次序得到的低一阶行列式，记作 $M_{ij}$；约定0阶行列式为1。}
\GlossaryEntry{模}{Modulus}{复数 $z=a+bi$ 对应点到原点的距离，$|z|=\sqrt {a^2+b^2}$。}
\GlossaryEntry{多重线性型}{Multilinear Form}{实向量空间上的实数值多输入函数，对每个输入位置分别保持加法与数乘；固定其他位置才考察该位置的线性。}
\GlossaryEntry{多重线性}{Multilinearity}{多输入函数固定其他输入后，对每一个输入分别保持加法与数乘；不是要求所有输入同时改变时仍整体线性。}
\GlossaryLetter{N}
\GlossaryEntry{$n$ 维实向量}{n-Dimensional Real Vector}{由 $n$ 个实数按固定次序组成的数组，是 $\mathbb R^n$ 的元素。}
\GlossaryEntry{必要条件}{Necessary Condition}{某结论成立时必须满足的条件；条件成立本身未必能推出该结论。例如，线性映射必须把零向量送到零向量。}
\GlossaryEntry{负矩阵}{Negative Matrix}{将原矩阵各元素变号所得的矩阵，表示把原映射的每个输出取相反向量。}
\GlossaryEntry{非齐次线性方程组}{Nonhomogeneous Linear System}{至少有一个常数项不为0的线性方程组 $A\boldsymbol x=\boldsymbol b$。}
\GlossaryEntry{非齐次递推}{Nonhomogeneous Recurrence}{带不恒为零的给定外部输入的线性状态更新，如 $\boldsymbol x_{k+1}=F\boldsymbol x_k+\boldsymbol u_k$；其通解是一串特解加任意齐次解序列。}
\GlossaryEntry{非零解}{Nontrivial Solution}{分量不全为0的解向量；不要求每个分量都非零。}
\GlossaryEntry{正规方程}{Normal Equations}{最小二乘问题对应的线性方程组 $A^{\mathrm T}A\boldsymbol x=A^{\mathrm T}\boldsymbol b$；它始终有解，全部解构成原矩阵零空间的一次平移。}
\GlossaryEntry{法向量}{Normal Vector}{垂直于所讨论平面或其方向子空间的非零向量；底面不退化时，叉乘与代数余子式向量提供相应的构造。}
\GlossaryEntry{归一化}{Normalization}{选定度量基准的条件；行列式规定标准基有序组或单位矩阵的值为1，以固定整体尺度。}
\GlossaryEntry{复数的 $n$ 次根}{nth Root of a Complex Number}{满足 $w^n=z$ 的复数 $w$；非零复数有 $n$ 个互异的 $n$ 次根。}
\GlossaryEntry{零空间}{Null Space}{线性映射把它映为零向量的全部输入所成的子空间，记为 $\operatorname {null}T$；矩阵的零空间就是齐次方程组的解空间。}
\GlossaryEntry{零度}{Nullity}{线性映射或矩阵的零空间的维数；等于相应齐次方程组自由变量的个数。}
\GlossaryEntry{数值稳定性}{Numerical Stability}{描述算法控制有限精度误差影响的性质；选主元可避免过小主元造成过大的消元倍数，但不保证一切误差消失。}
\GlossaryLetter{O}
\GlossaryEntry{斜投影}{Oblique Projection}{沿与目标子空间不垂直的互补方向投到该子空间、并保持其上的向量不动的线性映射；本章以沿 $(1,1,1)^{\mathrm T}$ 投到 $xy$ 平面为例。}
\GlossaryEntry{一一对应}{One-to-One Correspondence}{两类对象之间互相唯一确定、没有遗漏的配对关系；固定标准基后，$\mathbb R^n$ 到 $\mathbb R^m$ 的线性映射与 $m\times n$ 矩阵一一对应。}
\GlossaryEntry{有序基}{Ordered Basis}{指定了基向量先后次序的一组基，坐标必须按该次序排列。}
\GlossaryEntry{有向面积}{Oriented Area}{按两条边的先后次序记录正负的平行四边形面积；平面标准定向中逆时针为正，共线时为0。}
\GlossaryEntry{有向体积}{Oriented Volume}{按有序边的定向记录正负的体积；三维中由混合积给出，通常体积为其绝对值。}
\GlossaryEntry{正交补向量}{Orthogonal Complement Vector}{本章指按长度和定向选定的垂直向量，区别于由所有垂直向量组成的正交补子空间；二维标准正交基下取逆时针转90度的像，高维由有向底面体积与定向确定。}
\GlossaryEntry{正交互补子空间}{Orthogonal Complementary Subspaces}{在给定实向量空间中互相正交，且直和等于整个空间的两个子空间；每个向量均有唯一分解，二者维数之和等于环境空间的维数，并不要求维数相等。}
\GlossaryEntry{正交直和}{Orthogonal Direct Sum}{互相正交的子空间所构成的直和；其中每个向量的分量分解唯一。}
\GlossaryEntry{正交矩阵}{Orthogonal Matrix}{满足 $Q^{\mathrm T}Q=E$ 的实方阵；列向量单位且两两垂直，逆矩阵等于转置，保持标准点积；不一定表示旋转，也可能包含反射。}
\GlossaryEntry{正交投影}{Orthogonal Projection}{沿垂直于目标直线或平面的方向得到投影的对应规则；向 $xy$ 平面投影得到 $(x,y,0)^{\mathrm T}$，用平面坐标记录则为 $(x,y)^{\mathrm T}$。}
\GlossaryEntry{子空间正交}{Orthogonal Subspaces}{两个子空间中任取一支向量配对，标准点积均为零的关系；仅有正交并不保证二者之和为整个空间。}
\GlossaryEntry{正交性}{Orthogonality}{内积为0所刻画的关系，是几何垂直的推广；第7章展开。}
\GlossaryEntry{外积}{Outer Product}{列向量与行向量的矩阵乘积 $\boldsymbol u\boldsymbol v^{\mathrm T}$；第 $(i,j)$ 个元素为 $u_iv_j$，两向量均非零时秩为1。}
\GlossaryLetter{P}
\GlossaryEntry{平行四边形定则}{Parallelogram Rule}{以同一起点的两个向量为邻边作平行四边形，其从共同起点出发的对角线表示向量和。}
\GlossaryEntry{平行多面体}{Parallelotope}{由同一空间中的向量作为邻边、系数分别在0到1之间变化而得到的区域；边向量相关时退化到较低维子空间。}
\GlossaryEntry{偏导数}{Partial Derivative}{固定其余输入变量，只改变一个变量时的瞬时变化率；在二元曲面上对应固定另一坐标的截线切线斜率。}
\GlossaryEntry{部分选主元}{Partial Pivoting}{每一步在当前列尚未处理的行中选绝对值最大的元素，交换到主元行后再消元。}
\GlossaryEntry{特解}{Particular Solution}{一个确实满足给定方程组的具体解；一个特解未必代表全部解。}
\GlossaryEntry{被动旋转}{Passive Rotation}{向量不动，转动描述它的坐标架；新坐标由对应主动旋转矩阵的逆矩阵换算。}
\GlossaryEntry{排列}{Permutation}{把 $1,\ldots ,n$ 的全部下标按某一顺序重排，或相应的下标一一对应；行列式排列式中每行每列恰取一个元素。}
\GlossaryEntry{置换矩阵}{Permutation Matrix}{单位矩阵的行重新排列所得矩阵，每行每列恰有一个1；左乘它可重排行，且 $P^{-1}=P^{\mathrm T}$。}
\GlossaryEntry{状态平面}{Phase Plane}{用二维状态的两个分量作为坐标形成的平面；点代表状态，曲线代表状态演化，不一定是物体运动的实际空间。}
\GlossaryEntry{主元}{Pivot}{阶梯形非零行的首个非零元素，用于组织消元与确定主元变量。}
\GlossaryEntry{主元列}{Pivot Column}{阶梯形中含主元的列；原矩阵对应列号的列可用于选取极大无关组。}
\GlossaryEntry{主元位置}{Pivot Position}{阶梯形中主元所在的位置；讨论原矩阵时指与化简后主元对应的位置。}
\GlossaryEntry{像素}{Pixel}{数字图像中一个采样位置及其数值记录；灰度图像的像素值表示亮度，彩色图像通常记录多个颜色分量。}
\GlossaryEntry{像素基}{Pixel Basis}{每个基图案仅在一个像素位置取1、其他位置取0；按固定顺序展开后就是图像向量空间的标准基。}
\GlossaryEntry{平面向量}{Plane Vector}{位于一个平面内的向量，选定平面基后可用两个坐标表示。}
\GlossaryEntry{PLU分解}{PLU Factorization}{把换行记录到置换矩阵中的LU分解，本书约定写成 $PA=LU$；第3.8节为选读。}
\GlossaryEntry{极坐标}{Polar Coordinates}{用到原点的距离 $r$ 与方向角 $\theta $ 描述平面位置，满足 $x=r\cos \theta ,y=r\sin \theta $；面积微元为 $r\,\mathrm dr\,\mathrm d\theta $。}
\GlossaryEntry{极坐标形式}{Polar Form}{用模 $r$ 和辐角 $\theta $ 表示非零复数的形式 $r(\cos \theta +i\sin \theta )$，也称极式。}
\GlossaryEntry{多项式空间}{Polynomial Space}{以多项式为元素、按通常加法和数乘运算组成的向量空间；$P_n$ 表示次数不超过 $n$ 的情形。}
\GlossaryEntry{势函数}{Potential Function}{其梯度给出某个向量场的标量函数；物理中的保守力也常以负梯度表示。无旋场能否在整个区域具有势函数还取决于区域条件。}
\GlossaryEntry{原像}{Preimage}{若 $T(\boldsymbol x)=\boldsymbol y$，则输入 $\boldsymbol x$ 是输出 $\boldsymbol y$ 的一个原像；双射中每个目标向量有唯一原像。}
\GlossaryEntry{辐角主值}{Principal Argument}{按约定区间选出的唯一辐角，本附录采用 $(-\pi ,\pi ]$，记为 $\operatorname {Arg}z$。}
\GlossaryEntry{纯虚数}{Purely Imaginary Number}{实部为0且虚部非零的复数 $bi$。}
\GlossaryLetter{Q}
\GlossaryEntry{二次型}{Quadratic Form}{若干变量的齐次二次多项式；第10章研究其代数表示与几何形状。}
\GlossaryEntry{量化}{Quantization}{把连续或精细的数值归并为有限精度的代表值；本例通过除以指定间隔并取最近整数实现，会丢失精度，一般不是线性映射。}
\GlossaryEntry{查询}{Query}{注意力计算中用于与键比较的向量，常将各位置的查询按行组成矩阵 $Q$。}
\GlossaryLetter{R}
\GlossaryEntry{值域}{Range}{线性映射全部输出向量组成的目标空间的子空间，记为 $\operatorname {range}T$；矩阵的值域等于其列空间。}
\GlossaryEntry{列秩分解}{Rank Factorization}{也称CR分解，将秩为 $r$ 的矩阵写成 $A=CR$；$C$ 有 $r$ 个独立列，$R$ 有 $r$ 个独立行，分别记录核心列与组合系数。}
\GlossaryEntry{矩阵的秩}{Rank of a Matrix}{矩阵列向量组的秩，等于列向量张成空间的维数，也等于行阶梯形的主元数。}
\GlossaryEntry{向量组的秩}{Rank of a Set of Vectors}{向量组的任意极大线性无关组所含向量个数，等于其张成空间的维数。}
\GlossaryEntry{秩与零度定理}{Rank-Nullity Theorem}{有限维线性映射的定义域维数等于零空间维数与值域维数之和；对 $m\times n$ 矩阵，有 $\dim \operatorname {null}A=n-\operatorname {rank}(A)$。}
\GlossaryEntry{秩一矩阵}{Rank-One Matrix}{列向量组的秩为1的非零矩阵；所有列均可由同一个非零列向量数乘得到。}
\GlossaryEntry{实轴}{Real Axis}{复平面中虚部为0的点组成的坐标轴。}
\GlossaryEntry{实部}{Real Part}{复数 $a+bi$ 中的实数 $a$。}
\GlossaryEntry{实向量空间}{Real Vector Space}{以实数作为数乘系数的向量空间。}
\GlossaryEntry{递推关系}{Recurrence Relation}{用已有项或当前状态确定后续项或状态的关系；线性状态更新可写为 $\boldsymbol x_{k+1}=F\boldsymbol x_k+\boldsymbol u_k$。}
\GlossaryEntry{行最简形矩阵}{Reduced Row Echelon Form}{主元为1且主元列其他元素均为0的行阶梯形矩阵，也称简化行阶梯形矩阵。}
\GlossaryEntry{反射变换}{Reflection}{关于过原点直线或平面作镜像的线性变换；反射两次恢复原向量。}
\GlossaryEntry{反身性}{Reflexivity}{某种关系对每个对象与其自身都成立；如向量组与自身等价。}
\GlossaryEntry{正则化}{Regularization}{在拟合目标中加入对参数等对象的约束或惩罚，以改善模型的稳定性与泛化；本章示例采用参数平方和惩罚。}
\GlossaryEntry{余项}{Remainder}{原函数与截断多项式之间的差，如 $R_n(t)=f(t)-T_n(t)$。}
\GlossaryEntry{残差向量}{Residual Vector}{目标与实际输出之差 $\boldsymbol b-A\boldsymbol x$；最小二乘解的残差属于 $\operatorname {Null}(A^{\mathrm T})$，与列空间正交。}
\GlossaryEntry{黎曼和}{Riemann Sum}{把区域分成有限小块，将每块中的函数采样值乘以该块大小后求和；网格逐渐细分时用于描述积分的累积思想。}
\GlossaryEntry{右手定则}{Right-Hand Rule}{右手拇指指向转轴正向，四指弯曲方向规定为正向旋转；叉乘中四指从第一向量弯向第二向量，拇指指向法线，二者沿用同一三维正定向。}
\GlossaryEntry{单位根}{Roots of Unity}{方程 $w^n=1$ 的解，均在单位圆上，辐角相邻间隔为 $2\pi /n$。}
\GlossaryEntry{旋转}{Rotation}{绕原点转过固定有向角的变换；复平面上乘 $e^{i\theta }$ 表示旋转 $\theta $。}
\GlossaryEntry{旋转矩阵}{Rotation Matrix}{表示绕原点旋转的矩阵；平面逆时针旋转 $\theta $ 时，两列依次为 $(\cos \theta ,\sin \theta )^{\mathrm T}$ 与 $(-\sin \theta ,\cos \theta )^{\mathrm T}$。}
\GlossaryEntry{舍入误差}{Rounding Error}{计算或存储时只保留有限位数字，舍去或舍入其余数字所产生的偏差。}
\GlossaryEntry{行阶梯形矩阵}{Row Echelon Form}{零行置底、每个非零行的首个非零位置逐行右移的矩阵。}
\GlossaryEntry{行阶梯形方程组}{Row Echelon System}{方程按固定未知数次序排列，全零方程置底，其余方程的首个非零项逐行右移的形式。}
\GlossaryEntry{行等价}{Row Equivalence}{两个矩阵可通过有限次初等行变换互相得到的关系。}
\GlossaryEntry{行秩}{Row Rank}{矩阵行向量张成空间的维数；等于列秩。}
\GlossaryEntry{行空间}{Row Space}{矩阵各行竖排后张成的空间，等于转置矩阵的列空间；$m\times n$ 矩阵的行空间是 $\mathbb R^n$ 的子空间。第3.5节从转置引入。}
\GlossaryEntry{行向量}{Row Vector}{将分量按固定次序横排得到的向量形式，是相应列向量的转置。}
\GlossaryLetter{S}
\GlossaryEntry{鞍点}{Saddle Point}{有的方向靠近平衡点、有的方向离开的矢量场形态；本节通过一个分量增长、另一个分量衰减的例子观察。}
\GlossaryEntry{萨鲁斯法则}{Sarrus' Rule}{三阶行列式的速算规则：抄写前两列，将右下方向的三条对角线乘积相加、右上方向的三条对角线乘积相减；它重排三阶展开的六项，不适用于四阶及更高阶行列式。}
\GlossaryEntry{标量}{Scalar}{来自所选域、用于数乘向量的数。}
\GlossaryEntry{标量场}{Scalar Field}{给空间中每个位置指定一个数的函数，如高度、温度；输入可用位置向量表示，输出是标量。}
\GlossaryEntry{数乘}{Scalar Multiplication}{以一个标量乘一个向量；在 $\mathbb R^n$ 中，每个分量都乘同一个数。}
\GlossaryEntry{映射的数乘}{Scalar Multiplication of a Map}{把映射的每个输出同乘标量，$(\lambda T)(\boldsymbol x)=\lambda T(\boldsymbol x)$。}
\GlossaryEntry{有符号投影长度}{Scalar Projection}{向量 $\boldsymbol y$ 沿单位方向 $\boldsymbol q$ 的带符号分量 $\boldsymbol q\cdot \boldsymbol y=|\boldsymbol y|\cos \theta $；反向一侧为负，原始非单位探测器的点积还需乘探测器的长度。}
\GlossaryEntry{混合积}{Scalar Triple Product}{三维中的 $\boldsymbol w\cdot (\boldsymbol u\times \boldsymbol v)$；给出三条有序边围出的有向体积，循环换序不变、交换两条边变号。}
\GlossaryEntry{缩放变换}{Scaling Transformation}{将各坐标分量分别乘以固定实数的线性变换；所有缩放因子非零时可逆，因子为零时会压缩相应方向。}
\GlossaryEntry{剪切变换}{Shear Transformation}{保持一组平行方向不变、使另一方向按位置成比例偏移的线性变换；平面例为 $(x,y)\mapsto (x+ky,y)$。}
\GlossaryEntry{Sigmoid函数}{Sigmoid Function}{本书指函数 $\sigma (t)=1/(1+e^{-t})$，把实数映入0与1之间；数值处于该区间不自动保证概率预测准确。}
\GlossaryEntry{相似矩阵}{Similar Matrices}{存在可逆矩阵 $P$ 使 $B=P^{-1}AP$ 的两个同阶方阵，表示同一线性变换在不同基下的矩阵。}
\GlossaryEntry{相似变换}{Similarity Transformation}{用可逆矩阵建立 $B=P^{-1}AP$ 的关系，表示同一线性变换在不同基下的矩阵；第4章选读引入。}
\GlossaryEntry{奇异矩阵}{Singular Matrix}{不存在双侧逆的方阵。}
\GlossaryEntry{奇异值分解}{Singular Value Decomposition}{用两组正交方向和非负缩放量分解矩阵的方法；第11章选读。}
\GlossaryEntry{反对称矩阵}{Skew-Symmetric Matrix}{满足 $A^{\mathrm T}=-A$ 的实方阵，主对角线元素全为0。}
\GlossaryEntry{Softmax归一化}{Softmax}{对一组实数分别取指数，再除以所有指数之和，得到总和为1的正权重；注意力通常逐行进行该操作。}
\GlossaryEntry{解}{Solution}{同时满足方程组所有方程的一组有序数。}
\GlossaryEntry{解集}{Solution Set}{方程组全部解构成的集合。}
\GlossaryEntry{解空间}{Solution Space}{齐次线性方程组全部解组成的向量空间；非齐次方程组有解时，其解集一般是该空间的平移，不是子空间。}
\GlossaryEntry{连续函数空间}{Space of Continuous Functions}{区间上所有连续函数按逐点加法与数乘组成的向量空间，记为 $C[a,b]$。}
\GlossaryEntry{空间向量}{Space Vector}{日常三维空间中的向量，选定空间基后可用三个坐标表示。}
\GlossaryEntry{张成空间}{Span}{给定向量组的全部线性组合构成的集合，是包含该组向量的最小子空间。}
\GlossaryEntry{生成组}{Spanning Set}{能够张成给定空间的向量组，不要求线性无关。}
\GlossaryEntry{方阵}{Square Matrix}{行数与列数相同的矩阵；$n\times n$ 方阵称为 $n$ 阶方阵。}
\GlossaryEntry{稳定结点}{Stable Node}{附近状态向平衡点靠近而不绕行的典型线性矢量场形态；本节通过两个分量分别衰减的例子观察。}
\GlossaryEntry{标准基}{Standard Basis}{在 $\mathbb R^n$ 中由 $\boldsymbol e_1,\ldots ,\boldsymbol e_n$ 组成的有序基，第 $i$ 个向量仅第 $i$ 分量为1。}
\GlossaryEntry{状态向量}{State Vector}{把模型中足以描述当前状态、并结合更新规则与输入确定后续演化的量有序排成的向量；其分量与所选模型有关。}
\GlossaryEntry{状态转移矩阵}{State-Transition Matrix}{在线性离散模型中把当前状态映到下一状态的矩阵；本书以固定矩阵$F$和无外部输入的关系$\boldsymbol x_{k+1}=F\boldsymbol x_k$作初步说明。}
\GlossaryEntry{稳态响应}{Steady-State Response}{在本节周期受迫且齐次偏离衰减的模型中，长期保留下来的周期响应；稳态不意味着状态静止，一般特解不一定是稳态响应。}
\GlossaryEntry{刚度矩阵}{Stiffness Matrix}{在线性结构模型中把节点位移映射为平衡外力的矩阵；本章三节点教学模型写成 $K\boldsymbol u=\boldsymbol f$。}
\GlossaryEntry{流线}{Streamline}{在自治矢量场中处处与场向量相切的曲线；本书连续系统中用作状态轨线的几何描述。}
\GlossaryEntry{严格对角占优}{Strict Diagonal Dominance}{本书采用按行条件：每行主对角元绝对值严格大于其余元素绝对值之和；它是无换行LU消元可行的一类充分条件。}
\GlossaryEntry{子空间}{Subspace}{沿用母空间运算、对加法和数乘封闭的非空子集。}
\GlossaryEntry{代入消元法}{Substitution Method}{先由一个方程表示某未知数，再代入其他方程以减少未知数的方法。}
\GlossaryEntry{充分条件}{Sufficient Condition}{满足该条件就能推出相应结论的条件；与必要条件须加以区别。仅把零向量送到零向量不足以保证映射线性。}
\GlossaryEntry{子空间的和}{Sum of Subspaces}{所有 $\boldsymbol u+\boldsymbol v$ 构成的子空间，其中 $\boldsymbol u\in U$、$\boldsymbol v\in V$，记为 $U+V$。}
\GlossaryEntry{满射}{Surjection}{目标空间的每个元素都至少能由一个输入得到的映射。}
\GlossaryEntry{对称矩阵}{Symmetric Matrix}{满足 $A^{\mathrm T}=A$ 的实方阵，对称位置的元素相等。}
\GlossaryEntry{对称性}{Symmetry}{某种关系成立时交换两个对象，关系仍成立；如向量组等价关系。}
\GlossaryLetter{T}
\GlossaryEntry{切平面}{Tangent Plane}{平滑曲面在一点附近的一阶仿射近似，由该点高度和各输入方向的偏导数决定。}
\GlossaryEntry{泰勒多项式}{Taylor Polynomial}{用函数在某一点的有限阶导数作系数构造的多项式，本书用于选读的函数近似例子。}
\GlossaryEntry{词元}{Token}{分词器按规则将文本切分得到的基本处理单元，可为完整的词、词的一部分或标点等，不必与汉字或自然语言中的词一一对应。模型先将其映射为初始向量，再通过后续计算融入上下文信息。}
\GlossaryEntry{迹}{Trace}{方阵主对角线上元素的和，记作 $\operatorname {tr}A$；在线性向量场中表示体积的瞬时相对膨胀率。}
\GlossaryEntry{暂态响应}{Transient Response}{在本书选读的衰减模型中，由初始状态偏离特解所产生、并随时间消失的齐次响应；一般齐次偏离未必衰减。}
\GlossaryEntry{传递性}{Transitivity}{若对象一与二、二与三具有某种关系，则对象一与三也具有该关系。}
\GlossaryEntry{转置}{Transpose}{对偶映射在原标准基所对应的对偶基下的矩阵，满足 $\boldsymbol y\cdot (A\boldsymbol x)=(A^{\mathrm T}\boldsymbol y)\cdot \boldsymbol x$；行列互换是其坐标计算形式，通常不是逆矩阵。向量横排与竖排的转换也沿用上标 $\mathrm T$。}
\GlossaryEntry{三角形定则}{Triangle Rule}{把两个向量首尾相接，从第一支起点指向第二支终点的向量表示它们的和。}
\GlossaryEntry{三对角矩阵}{Tridiagonal Matrix}{除主对角线及其上、下两条次对角线以外的元素均为零的方阵；常见于相邻节点耦合的线性模型，其行列式可利用带状结构降阶递推。}
\GlossaryEntry{三重积分}{Triple Integral}{对空间区域上小体积的函数贡献进行累积的积分；如把体密度乘体积微元后相加得到总质量。}
\GlossaryEntry{零解}{Trivial Solution}{所有未知数均为0的解，是任意齐次线性方程组的一个解。}
\GlossaryEntry{平凡子空间}{Trivial Subspaces}{所讨论空间的零子空间与全空间；零空间情形二者重合。}
\GlossaryEntry{二维广义叉乘}{Two-Dimensional Generalized Cross Product}{本章采用一个平面向量输入的垂直生成操作，$(a,b)^{\mathrm T}\mapsto (-b,a)^{\mathrm T}$；区别于三维的双向量叉乘。}
\GlossaryEntry{双塔召回模型}{Two-Tower Retrieval Model}{分别生成用户与候选内容的向量表征，再按匹配程度筛选候选的模型；两座塔通常包含非线性运算。}
\GlossaryLetter{U}
\GlossaryEntry{唯一解}{Unique Solution}{方程组恰有一个解向量；有解且没有自由变量时出现。}
\GlossaryEntry{单位圆}{Unit Circle}{复平面上满足 $|z|=1$ 的点集，即以原点为圆心、半径为1的圆。}
\GlossaryEntry{单位下三角矩阵}{Unit Lower Triangular Matrix}{主对角线元素全为1的下三角矩阵，其对角线下方元素可为任意实数。}
\GlossaryEntry{单位向量}{Unit Vector}{长度为1的向量。}
\GlossaryEntry{上三角矩阵}{Upper Triangular Matrix}{主对角线下方全为零的方阵。}
\GlossaryEntry{用户—物品交互矩阵}{User-Item Interaction Matrix}{用行对应用户、列对应物品，按位置记录行为或反馈的矩阵；未记录到行为不自动等于不喜欢。}
\GlossaryLetter{V}
\GlossaryEntry{值}{Value}{注意力计算中被权重加权汇聚的信息向量，常按行组成矩阵 $V$。}
\GlossaryEntry{向量}{Vector}{向量空间中的元素；在本书主线中主要表现为有序实数组，低维时可画成箭头。}
\GlossaryEntry{向量加法}{Vector Addition}{在 $\mathbb R^n$ 中将相同位置的分量分别相加的运算。}
\GlossaryEntry{向量方程}{Vector Equation}{两端为向量表达式的等式；本章重点为 $x_1\boldsymbol a_1+\cdots +x_n\boldsymbol a_n=\boldsymbol b$。}
\GlossaryEntry{矢量场}{Vector Field}{为状态空间的每个点指定一支向量的对应；连续系统的 $A\boldsymbol y$ 在该点表示瞬时变化率。}
\GlossaryEntry{向量组}{Vector List}{按次序列出的有限个向量，允许出现重复向量和零向量。}
\GlossaryEntry{向量空间}{Vector Space}{配有封闭的加法与数乘并满足相应八条运算公理的非空集合。}
\GlossaryEntry{向量减法}{Vector Subtraction}{加上相反向量的运算，$\boldsymbol u-\boldsymbol v=\boldsymbol u+(-\boldsymbol v)$。}
\GlossaryEntry{黏性阻尼}{Viscous Damping}{采用阻力与速度成正比、方向相反的线性近似；在一维教学模型中写为 $-c\dot x$。}
\GlossaryEntry{体积缩放倍数}{Volume Scale Factor}{实方阵表示的线性变换对区域体积的乘数，等于行列式的绝对值；退化变换的倍数为0。}
\GlossaryLetter{W}
\GlossaryEntry{权重矩阵}{Weight Matrix}{包含模型可训练系数、用于线性映射的矩阵；注意力中的 $A$ 也称权重矩阵，但它由当前输入计算得到，与可训练投影参数 $W_Q,W_K,W_V$ 不同。}
\GlossaryEntry{全空间}{Whole Space}{所讨论的整个向量空间，它也是自身的子空间。}
\GlossaryLetter{Z}
\GlossaryEntry{零矩阵}{Zero Matrix}{所有元素都为0的矩阵。}
\GlossaryEntry{零子空间}{Zero Subspace}{只含零向量的子空间 $\{\boldsymbol 0\}$，其基可取空组，维数为0。}
\GlossaryEntry{零向量}{Zero Vector}{加法的零元素；在 $\mathbb R^n$ 中各分量均为0，几何上长度为0且没有确定方向。}
